\section{A complete scattered first factor and its failure to compose}
\label{sec:completion}

A source decoder does not automatically extend to a bijection on all
matrices.  For the scattered normal-basis source, such an extension does
exist.  The particular extension below nevertheless forces every receiver
of the remaining factor to have its own class.  Both assertions concern
the complete finite state space; neither supplies an efficient evaluator
for the orbit permutation.

\subsection{The source and a total permutation}

Let $d\ge2$ and $b\ge1$, put
$V=(\F_2^b)^{d\times d}$, and let
$B=(\F_2^b)^d$ be the whole-column alphabet, of cardinality
$q=2^{db}$.  Use the reversed normal-basis identification
$\F_2^d\cong E=\F_{2^d}$ underlying
\cref{thm:interval-attainment}; a column rotation by $s$ acts on a row
as $\Phi^s$, where $\Phi(x)=x^2$.
Let $\beta$ be its normal element.  Its $d$ conjugates are distinct,
and none equals $0$ or $1$.
For $u\in B$, write $a_\ell\in E$ for the field encoding of
$(u_{j,\ell})_{j\in\Z/d}$ and define $w(u)$ by its encoded rows:
\[
 w(u)_{0,\cdot,\ell}=a_\ell,\qquad
 w(u)_{i,\cdot,\ell}=\beta a_\ell\quad(i\ne0).
\]
Set $W=\{w(u):u\in B\}$.  Also write $\bar u$ for the matrix whose
every column is $u$, and let $z(u)$ have column zero equal to zero and
all other columns equal to $u$.

\begin{theorem}[A complete scattered first factor]
\label{thm:scattered-completion}
\formal{CompositionScatteredSource.lean}{scatteredSource_finrank,scattered_orbit_injective,scattered_punctured_orbits_disjoint}
\formal{CompositionScatteredCompletion.lean}{scatteredCompletion_equivariant,scatteredFirstFactor_classUniform,scatteredFirstFactor_source,scatteredFirstFactor_fixed_preimage}
\formal{CompositionSparsePatch.lean}{columnPatch_involutive,columnPatch_nonzero_column}
The binary-linear source $W$ has dimension $db$.  There are total
permutations $H,R,K_1$ of $V$ such that $H$ commutes with column
rotation, $R$ changes only column zero, and $K_1=RH$ is class-uniform
for the two classes $\{0\}$ and $(\Z/d)\setminus\{0\}$.
They satisfy
\[
 H(w(u))=z(u),\qquad K_1(w(u))=\bar u\quad(u\in B),\qquad
 K_1^{-1}(\operatorname{Fix}(T))=W.
\]
Here $R$ is an involution defined on every input.
\end{theorem}
\begin{outline}
The row-zero projection recovers $u$, so $w$ is an injective binary-linear
map and $\dim W=db$.  If
$T^sw(u)=T^tw(v)$ with $u,v\ne0$, choose a nonzero lane parameter.
The two row types give
\[
 \Phi^s(a)=\Phi^t(a')\ne0,\qquad
 \Phi^s(\beta)\Phi^s(a)=\Phi^t(\beta)\Phi^t(a').
\]
Cancellation and the distinct conjugates of $\beta$ give $s=t$, then
all lane parameters agree.  Thus $(s,u)\mapsto T^sw(u)$ is injective
for $u\ne0$.  Similarly $(s,u)\mapsto T^sz(u)$ is injective: its
unique zero column determines $s$, and any other column determines $u$.

These two orbit families are disjoint.  Suppose a rotated $w(u)$ equals
a rotated $z(v)$.  Choose a lane whose row-zero field parameter is
nonzero.  On the punctured side, a row in this lane is either zero or
the common punctured binary word, according as the corresponding bit of
$v$ is zero or one.  Both rows zero and one on the scattered side are
nonzero.  They must therefore be equal on the punctured side.  In field
coordinates this would give
$\Phi^s(\beta)\Phi^s(a)=\Phi^s(a)$ with $a\ne0$, contradicting
$\Phi^s(\beta)\ne1$.

Define $H$ to exchange $T^sw(u)$ and $T^sz(u)$ for every $s$ and
every $u\ne0$, and fix every remaining state.  Injectivity and
disjointness make these disjoint pairs, so this is an actual total
involution.  Rotation permutes the pairs with their orientation
preserved, hence $HT=TH$.  Zero lies outside both families and is fixed.
This specializes the completion to the present source and needs no
assumed equivariant extension.

Define $R$ to exchange $z(u)$ and $\bar u$ for each $u\ne0$, fixing
all other inputs.  Each pair agrees outside column zero, so $R$ leaves
those columns unchanged.  The pairs are disjoint: a punctured state has
zero column zero, whereas a nonzero constant state does not.  Thus $R$
is again a total involution and $Rz(u)=\bar u$, including $u=0$.

For the exceptional class use the decoder
$Y\mapsto\operatorname{col}_0(RH(Y))$; for the other class use
$Y\mapsto\operatorname{col}_0H(Y)$.
If $c\ne0$, equivariance gives
$\operatorname{col}_0H(T^cX)=\operatorname{col}_cH(X)
 =\operatorname{col}_cK_1(X)$.
The exceptional decoder gives column zero by definition.  This proves
sharedness on every input.  Finally the $q$ matrices $\bar u$ exhaust
$\operatorname{Fix}(T)$, and bijectivity of $K_1$ gives the stated
preimage equality.
\end{outline}

\begin{corollary}[The exceptional class is necessary]
\label{cor:scattered-one-class}
\formal{CompositionScatteredCompletion.lean}{scattered_one_class_impossible}
No injective equivariant map sends all of $W$ into
$\operatorname{Fix}(T)$.
\end{corollary}
\begin{outline}
Choose $u\ne0$. Equivariance and the fixed image imply
\[K(Tw(u))=TK(w(u))=K(w(u)).\]  Injectivity forces $Tw(u)=w(u)$,
contradicting the free-orbit calculation above.  A one-class complete
map is equivariant, so it is excluded as well.
\end{outline}

\subsection{Complete-input disagreement}

The agreement estimates in \cref{prop:agreement} have the following
robust form.  The comparison need not be a permutation.

\begin{proposition}[Robust comparison with an equivariant map]
\label{prop:robust-equivariant-comparison}
\formal{CompositionRobustAgreement.lean}{agreement_error_cover,robust_full_input_count,robust_repeated_class_count}
Let $A$ be any finite alphabet with at least two elements and put
$q=|A|^d$.  Let $K_1,H\colon A^{d\times d}\to A^{d\times d}$ be
total maps, with $H$ equivariant, and let $K_2$ be class-uniform for a
partition $p$ with $M$ nonempty classes of sizes $n_J$.  Under the
uniform law on the complete input space, put
\[
 \gamma=\Pr[K_2K_1(X)\ne\tau X],\qquad
 \epsilon=\Pr[K_1(X)\ne H(X)].
\]
Then
\[
 \gamma+\epsilon\ \ge\
 1-\frac Md-\frac{\sum_Jn_J(n_J-1)}{2dq}.
\]
If a receiver class is repeated, also
$\gamma+\epsilon\ge(1-q^{-1})/2$.
\end{proposition}
\begin{outline}
The map $K_2H$ has the same receiver partition, with decoder
$f_J\circ H$.  If $K_2K_1$ is correct at $X$ and $K_1(X)=H(X)$,
then $K_2H$ is correct there.  The entire input space is therefore
covered by the correct inputs of $K_2H$, the failures of $K_2K_1$,
and the disagreements between $K_1$ and $H$.  Sum their cardinalities
and apply the two bounds of \cref{prop:agreement} to $K_2H$.
The formal statements retain the integer counts and multiply through
by the positive denominators.  They use no injectivity of $H$ or $K_1$.
\end{outline}

\begin{theorem}[Sparse completion forces a fully separated second factor]
\label{thm:scattered-second-factor}
\formal{CompositionSparsePatch.lean}{columnPatch_support_card}
\formal{CompositionSparseObstruction.lean}{scatteredFirstFactor_disagreement,scattered_remaining_factor_singletons}
For the binary construction of \cref{thm:scattered-completion},
\[
 |\{X:R(X)\ne X\}|=2(q-1),\qquad
 \Pr[K_1(X)\ne H(X)]=\frac{2(q-1)}{q^d}.
\]
If $d\ge3$ and a class-uniform total map $K_2$ satisfies
$K_2K_1=\tau$ on all inputs, every receiver class of $K_2$ is a
singleton.  Thus its uniquely required remaining factor
$\tau K_1^{-1}$ needs $d$ classes.
\end{theorem}
\begin{outline}
The support of $R$ consists exactly of the two states $z(u),\bar u$
for every $u\ne0$.  Bijectivity of $H$ transports this support count
to the disagreement count for $RH$ and $H$.
If $K_2$ had a repeated class, the previous proposition with
$\gamma=0$ would force
\[
 \frac{2(q-1)}{q^d}\ge\frac{q-1}{2q},
 \qquad\text{hence}\qquad q^{d-1}\le4.
\]
But $q=2^{db}\ge8$ when $d\ge3$ and $b\ge1$.
The formal proof uses the equivalent integer inequalities, avoiding
subtraction or division of probabilities.  No exclusion at $d=2$ is
asserted.
\end{outline}

This disagreement counts complete input states.  It is different from
the number of output columns modified by a map, a linear error rank,
or a time charge.  A change confined to one output column can affect
almost every input.

\subsection{Endpoint classes are not multiplicative}

The need to retain actual factors is already visible in a short
bijective control.  For binary matrices define rowwise maps
\[
 (B_1X)_{i,c}=\begin{cases}X_{i,0}+X_{i,1},&c=0,\\X_{i,c},&c\ne0,\end{cases}
 \qquad
 (B_2X)_{i,c}=X_{i,c}+\sum_jX_{i,j}.
\]

\begin{proposition}[Two classes followed by one can require all classes]
\label{prop:endpoint-class-blowup}
\formal{CompositionEndpointBlowup.lean}{endpointFirst_involutive,endpointSecond_involutive,endpointFirst_classUniform,endpointSecond_equivariant,endpoint_composite_getter,endpoint_composite_requires_singletons}
For every even $d\ge2$, the maps $B_1,B_2$ are total involutions.
The first uses the singleton/rest partition and the second is equivariant.
Nevertheless every class-uniform representation of $B_2B_1$ has $d$
singleton classes.
\end{proposition}
\begin{outline}
The first map adds the unchanged column one twice when squared.
For the second, let $s_i=\sum_jX_{i,j}$.  The row sum after the map
is $s_i+d s_i=s_i$ because $d$ is even in $\F_2$, so its square is
also the identity.  The sum is rotation-invariant, proving that $B_2$
is equivariant; the two stated decoders for $B_1$ are
$Y\mapsto Y_{\cdot,0}+Y_{\cdot,1}$ and $Y\mapsto Y_{\cdot,0}$.

The normalized getter of the composite at receiver $c$ is
\[
 g_c(Y)_i=Y_{i,0}+\sum_jY_{i,j}
   +\begin{cases}0,&c=0,\\Y_{i,1-c},&c\ne0.\end{cases}
\]
These functions are pairwise distinct.  Comparing $c=0$ with $e\ne0$,
use a row supported at $1-e$.  Comparing nonzero distinct $c,e$, use
a row supported at $1-c$; the extra term at $1-e$ vanishes.
A shared class would make the corresponding normalized getters equal
on every input, a contradiction.  The formal proof derives this getter
equality from the complete class-uniform interface.
\end{outline}
